% Section 8 of the manuscript: text recoverability (question 4). The method is stated first, as the outline
% requires; it follows sections 6.8 and 10.7 of the frozen registration. The results are written from the result
% files e2_closed_set_<cell>.json and e2_open_set_<cell>.json after the analysis.

\section{Results: text recoverability}
\label{sec:recoverability}

\subsection{Method}
\label{sec:recoverability-method}

The text-recoverability study asked how much of the typed text the key rows of a published log still reveal. Two
attacks were run on the observation log of each cell separately, with the plan, the procedures and the settings that
built it; the attacks refused a log whose plan or settings did not match. Both read the category and the time of
every key row of a press and nothing else: they did not use the session rows or the rows of key releases, and no plan
sent a modifier key, so on Windows, where modifier keys produce rows, no such row was read either. The ground truth of
an attack was the reading that was typed,
not the original sentence. The trials of scenario S12, whose phrase is a generated kana string, were not scored. When
every block is scored this leaves 2{,}925 trials with live conversion on and 2{,}700 with it off on macOS, and 2{,}700
on Windows.

\paragraph{Assigning presses to trials.} Both attacks assign the observed presses to the trials in plan order, block
by block. A block is placed where its marker of four arrow presses and all its trials match the observed presses. A
block whose observed presses do not agree with the plan (a press missing, added or changed) is not scored: its trials
are counted as unscored, and the assignment continues with the next block that matches, so that a difference loses
only the block that holds it. This behaviour was checked by the self-test of the code on the cases it constructs, not
proved in general. Before each freeze the assignment, without any attack, was run on the key rows of the pilots: it
scored every block of the three macOS pilot runs that used the arrow-key marker (15 blocks, 282 trials) and of the
seven Windows pilot runs (27 blocks, 429 trials) and left no press unassigned. These logs had no missing or added
press, so the check shows that clean logs are scored whole; it does not exercise the recovery after a difference.
Accuracies and error rates are computed over the scored trials, and the numbers of unscored blocks and trials are
reported per cell.

\paragraph{Closed set.} For each trial the attacker was given $K$ candidate phrases, the true one and $K-1$ decoys
drawn from the phrases that the trial's scenario could have produced (excluding phrases with the same reading). It
knew the scenario, the system, the speed, the generator, the procedures and the settings, generated the expected log
of each candidate with the same generator, and chose the candidate whose expected log was closest to the observed
one. $K$ was 10, 100 and 1{,}000. The attack was repeated on six feature sets that remove information from the log,
named here, in the figures and in the tables as in the registration: exact times (\texttt{all}), no times
(\texttt{no\_time}), times in 1-second bins (\texttt{bin1s}), and each of these three with every run of letter keys
made one token (\texttt{no\_runs}, \texttt{bin1s\_count} and \texttt{counts\_only}). The registration describes the
last three as having no lengths of letter runs, but the token of a run kept the number of its presses, and the distance
charged the difference of these numbers; what the three sets removed was the time of every press of a run after its
first (deviation D13). Without times, \texttt{counts\_only} therefore held the same information as \texttt{no\_time},
the category of every press in order, and differed from it only in how the distance counted a run of letter keys.
Because an attacker given a detailed log can always discard
information, the accuracy of a feature set is the maximum over that set and every coarser set derived from it. The
accuracy is reported with its Wilson interval and, descriptively, with a 95\% percentile interval from 2{,}000
resamples of whole blocks, together with the chance level $1/K$, the mean rank of the true phrase and, without an
interval, the accuracy of each feature set alone per scenario and per speed.

\paragraph{The near-chance rule.} The synthetic key intervals were drawn independently of the keys, so the logs carry
none of the timing structure of human typing, neither the timing of particular key pairs nor the pauses between words
and phrases. The accuracy on the four feature sets that keep any time is therefore reported as an estimate only and
does not enter the recommendation, which considers the two sets that keep no time (\texttt{no\_time} and
\texttt{counts\_only}, registered as its coarsening). The accuracy of a feature set is near chance at $K=10$ or $K=100$
when the upper
end of its 95\% Wilson interval is at most $2/K$. A set is near chance for the recommendation when it is near chance at
both values of $K$ in both macOS cells, and the recommended granularity is the more detailed of the two sets that is
near chance in this sense; if neither is, no granularity is recommended, and this is reported. At $K=1{,}000$ an attack
at the chance level would rarely be judged near chance at the planned numbers of trials, so $K=1{,}000$ is reported
descriptively and is not part of the rule. The probability of the judgement at the planned numbers of trials was fixed
in the registration and is given with the results (Table~\ref{tab:nearchance}). The margin of $2/K$ and the values of
$K$ are choices of the design, not derived from a criterion of privacy. The Wilson interval treats trials as
independent; with positive dependence within a block the rule can be met more easily than it should be, which is why the
block-resampled interval is reported next to it. The recommendation was decided from the two macOS cells. The same
judgement was made for the Windows cell and, where it does not hold, is reported as a qualification of the
recommendation for logs recorded on Windows; it does not change the recommendation.

\paragraph{Exploratory: feature sets without run lengths.} Because no registered set dropped the lengths of letter
runs, the closed-set attack was also run, as an exploratory analysis that is not registered (deviation D13), on two
feature sets without times that drop them: \texttt{letter\_runs\_unit}, in which every run of letter keys is one token
whatever its length and every other press one token, and \texttt{all\_runs\_unit}, in which every run of presses of the
same category is one token whatever its length, so that the numbers of Backspace, arrow and other repeated presses are
dropped as well. Both keep the order of the tokens. The predicted logs, the distance, the candidates and the
tie-breaking of each trial, the values of $K$ and the intervals were those of the registered attack, and every run also
computed \texttt{no\_time} and \texttt{counts\_only} and stopped unless they equalled the registered results. As
\texttt{all\_runs\_unit} can be derived from \texttt{letter\_runs\_unit}, the accuracy of \texttt{letter\_runs\_unit}
is the larger of the two. Both sets can also be derived from \texttt{no\_time} and \texttt{counts\_only}; at every $K$
and in every cell each was less accurate than both, so the registered accuracies, their Wilson intervals and the
registered judgement are unchanged. The registered near-chance rule was applied to both sets as an exploratory judgement, from
the two macOS cells, with the Windows judgement as its qualification; it is reported next to the registered judgement
and does not change the recommendation. The two sets, the rule and the place of the result in this article were fixed
after all registered results had been seen and before the analysis was run, and a related measurement on single trials
of the Windows cell, made with a related attack in another study, had been seen before (deviation D13).

\paragraph{Open set.} A second attacker saw only the categories of the key presses and, with a character language
model, searched for the kana string that Romaji input would spell with exactly the observed letter keys and the
observed punctuation, maximising its probability. Two levels of knowledge were evaluated, named K0 and K1 as in the
registration (the names are not the number $K$ of candidates of the closed set). K0 knew nothing of the genre beyond
the frequency of single characters (a character unigram model estimated from the training phrases). K1 knew the
genre but not the individual work: a character trigram model with Witten--Bell smoothing
\citep{wittenbell1991}, trained on the corpus with the whole works of the evaluation phrases held out (three of the
twelve works, drawn with the seed, holding at least a fifth of the phrases; a training phrase with the reading of an
evaluation phrase was removed). A stronger attacker of the second level, K1+, for which the corpus was split by phrase
instead of by work, is reported alongside it. The baseline was the same model without the log, drawing a string from
its prior while knowing either the number of segments or the true length. Reported are the character error rate of
the attack and of the baseline, their paired difference, and the difference from a negative control in which each
trial is decoded from another trial's observation, each with a percentile interval from resampling whole blocks. The
open set does not enter the rule; it is reported next to the recommendation as a second view of what the lengths of
letter runs, which the \texttt{no\_time} set keeps, reveal.

\paragraph{What is not evaluated.} The rule judges the identification of one phrase among $K$ candidates in one
trial; it is not a criterion of privacy. Combining several trials of one writer, matching a log against known final
texts, and identifying a writer from the timing of a log are not evaluated, and no log is described as safe to
publish. The session rows carry their
times and durations in milliseconds and could undo the coarsening of the key rows; they were not given to the attacks,
and publishing them is not evaluated. Neither are the rows of key releases or, on Windows, the rows of modifier keys.
The feature sets that keep time are reported as estimates only, so whether key rows with times can be published is
not judged, and granularities between milliseconds and 1-second bins were not evaluated. Neither attack is
claimed to be the strongest possible. The attacker knew the generator and the procedures, and the input had no
typing errors; these conditions are expected to favour identification, which is stated as an expectation, not as a
result.

\subsection{Results}
\label{sec:recoverability-results}

Every block of every cell was scored (\rv{macos-standard-live-off.q4.unscored} unscored with live conversion off,
\rv{macos-standard-live-on.q4.unscored} with it on, \rv{windows-msime-tsf.q4.unscored} on Windows), so the rule was
applied to the planned numbers of trials, whose properties Table~\ref{tab:nearchance} gives.

\paragraph{Closed set and the recommendation.} Table~\ref{tab:closed-set} gives the accuracy of the two feature sets
that keep no time. With live conversion off, the accuracy of \texttt{no\_time} was
\rv{macos-standard-live-off.q4.no_time.k10.acc} at $K = 10$ and \rv{macos-standard-live-off.q4.no_time.k100.acc} at
$K = 100$; with it on, \rv{macos-standard-live-on.q4.no_time.k10.acc} and
\rv{macos-standard-live-on.q4.no_time.k100.acc}. The accuracy of \texttt{counts\_only}, the coarsest registered set,
from which no other registered set is derived, was \rv{macos-standard-live-off.q4.counts_only.k10.acc} and
\rv{macos-standard-live-off.q4.counts_only.k100.acc} with live conversion off and
\rv{macos-standard-live-on.q4.counts_only.k10.acc} and \rv{macos-standard-live-on.q4.counts_only.k100.acc} with it on.
In both macOS cells the upper end of every one of these Wilson intervals lies above $2/K$, so neither set is near chance
under the registered rule, which is not a criterion of privacy (Section~\ref{sec:recoverability-method};
\texttt{no\_time}: \rv{macos-standard-live-off.q4.no_time.near} and
\rv{macos-standard-live-on.q4.no_time.near}; \texttt{counts\_only}: \rv{macos-standard-live-off.q4.counts_only.near}
and \rv{macos-standard-live-on.q4.counts_only.near}), and the recommended granularity is \rv{q4.recommendation}: no
granularity of the key rows is recommended. In both macOS cells, at $K = 10$ and $K = 100$, the
maximum that gives the accuracy of \texttt{no\_time} came from the attack on \texttt{counts\_only}, which holds the
same information as \texttt{no\_time} and differs from it only in how its distance counts a run of letter keys
(Section~\ref{sec:recoverability-method}): counted as one token of its length, the runs identified the phrase at least
as often as the attack on \texttt{no\_time} alone did (supplement). The decoys were drawn from the phrases of the
trial's scenario, whose readings have 6 to 30 kana (Section~\ref{sec:method-plans}), and were not matched to the true
phrase in length.

On Windows, under the same judgement, \rv{q4.windows.qualification}. The Windows accuracies of the two sets equal those
of the macOS cell with live conversion off to the reported precision. The two plans were generated from the same seed
and share every trial's block, scenario, speed, phrase and times; their presses differ only in S2, where the macOS
procedure presses one Return more, and in S4, where the Windows procedure presses one Escape more (deviation D1), and
the presses of letters, digits and punctuation are the same in every trial. The Windows judgement is therefore not an
independent replication of the macOS one for these sets.

% Every word of the exploratory judgement is generated from the result files (q4x.*); the paragraph was to be here
% whatever the result (deviation D13).
\paragraph{Exploratory: feature sets without run lengths.} In the exploratory analysis of deviation D13, the accuracy of
\texttt{letter\_runs\_unit} was \rv{macos-standard-live-off.q4x.letter_runs_unit.k10.acc} at $K = 10$ and
\rv{macos-standard-live-off.q4x.letter_runs_unit.k100.acc} at $K = 100$ with live conversion off and
\rv{macos-standard-live-on.q4x.letter_runs_unit.k10.acc} and \rv{macos-standard-live-on.q4x.letter_runs_unit.k100.acc}
with it on; that of \texttt{all\_runs\_unit}, which also drops the numbers of repeated presses of the other categories,
was \rv{macos-standard-live-off.q4x.all_runs_unit.k10.acc} and \rv{macos-standard-live-off.q4x.all_runs_unit.k100.acc}
with live conversion off and \rv{macos-standard-live-on.q4x.all_runs_unit.k10.acc} and
\rv{macos-standard-live-on.q4x.all_runs_unit.k100.acc} with it on. Under the registered rule, applied as an exploratory
judgement, \rv{q4x.judgement}: in both cells, \rv{q4x.byk} (\texttt{letter\_runs\_unit} with live conversion off
and on: \rv{macos-standard-live-off.q4x.letter_runs_unit.k10.near} and
\rv{macos-standard-live-on.q4x.letter_runs_unit.k10.near} at $K = 10$,
\rv{macos-standard-live-off.q4x.letter_runs_unit.k100.near} and
\rv{macos-standard-live-on.q4x.letter_runs_unit.k100.near} at $K = 100$; \texttt{all\_runs\_unit}:
\rv{macos-standard-live-off.q4x.all_runs_unit.k10.near} and \rv{macos-standard-live-on.q4x.all_runs_unit.k10.near},
\rv{macos-standard-live-off.q4x.all_runs_unit.k100.near} and \rv{macos-standard-live-on.q4x.all_runs_unit.k100.near}).
Applied to these two sets, the rule would recommend \rv{q4x.granularity}. On Windows, under the same judgement,
\rv{q4x.windows}; as for the registered sets, its values equal those of the macOS cell with live conversion off. The
registered judgement and recommendation do not change. The supplement gives both sets at every $K$, with the accuracy
per speed and per scenario.

Figure~\ref{fig:closed-set} shows every feature set at every $K$, and the supplement gives their values, with the
accuracy per speed and scenario. The sets that keep time are estimates only and enter neither the judgement nor the
recommendation (Section~\ref{sec:recoverability-method}). The accuracy of an attack on a set alone was not monotone in
the information the set keeps: at $K = 10$ the attack on \texttt{all} alone identified fewer phrases than the attack on
\texttt{counts\_only} in both macOS cells (supplement), so the attacks are not the strongest possible, and the maximum
over the coarser sets is what the figure shows.

\paragraph{Open set.} Table~\ref{tab:open-set} gives the character error rate of the attack and of the prior-only
baseline, their difference (gain), and the gain over the shuffled control, in which each trial is decoded from another
trial's observation. With the number of segments known, the gain of K1 was
\rv{macos-standard-live-off.q4.open.K1.count.gain} with live conversion off and
\rv{macos-standard-live-on.q4.open.K1.count.gain} with it on, and its gain over the shuffled control
\rv{macos-standard-live-off.q4.open.K1.count.gainvsshuffled} and
\rv{macos-standard-live-on.q4.open.K1.count.gainvsshuffled}; for K1+ the gain over the shuffled control was
\rv{macos-standard-live-off.q4.open.K1plus.count.gainvsshuffled} and
\rv{macos-standard-live-on.q4.open.K1plus.count.gainvsshuffled}, and for K0
\rv{macos-standard-live-off.q4.open.K0.count.gainvsshuffled} and
\rv{macos-standard-live-on.q4.open.K0.count.gainvsshuffled}. With the true length known, the gain of K1 over the
baseline was \rv{macos-standard-live-off.q4.open.K1.true.gain} and \rv{macos-standard-live-on.q4.open.K1.true.gain}. The
open-set file of the Windows cell is byte-identical to that of the macOS cell with live conversion off: the attack
decodes the letter keys and the punctuation, which the two plans share, and every press other than the input-mode keys
reached the log in both cells. The open set does not enter the rule; it is a second view of what the letter-run lengths
kept by \texttt{no\_time} reveal, for these attacks.

% Generated by manuscript/tools/results_tables.py from the result files; do not edit by hand.
\begin{sidewaystable}[p]
\caption{Closed-set identification for the two feature sets that keep no time: the accuracy (the maximum over the set and the sets derivable from it) with its Wilson interval and, descriptively, the interval that resamples blocks; a set is near chance at K when the upper end of the Wilson interval is at most 2/K, and near chance for the rule when it is so at both K = 10 and K = 100. The recommendation is decided from the two macOS cells; the Windows judgement is reported as its qualification; --: no value}
\label{tab:closed-set}
\begin{tabular}{L{0.20}L{0.12}L{0.06}L{0.07}L{0.22}L{0.15}L{0.07}L{0.10}}
\toprule
Cell & Feature set & K & Trials & Accuracy & Block interval & 2/K & Near chance \\
\midrule
macOS, live conversion off & \texttt{no\_time} & 10 & 2{,}700 & 0.902 [0.890, 0.913] & 0.891, 0.913 & 0.200 & no \\
 & \texttt{no\_time} & 100 & 2{,}700 & 0.421 [0.403, 0.440] & 0.408, 0.442 & 0.020 & no \\
 & \texttt{counts\_only} & 10 & 2{,}700 & 0.902 [0.890, 0.913] & 0.890, 0.913 & 0.200 & no \\
 & \texttt{counts\_only} & 100 & 2{,}700 & 0.421 [0.403, 0.440] & 0.403, 0.440 & 0.020 & no \\
macOS, live conversion on & \texttt{no\_time} & 10 & 2{,}925 & 0.907 [0.896, 0.917] & 0.897, 0.917 & 0.200 & no \\
 & \texttt{no\_time} & 100 & 2{,}925 & 0.435 [0.417, 0.453] & 0.418, 0.454 & 0.020 & no \\
 & \texttt{counts\_only} & 10 & 2{,}925 & 0.907 [0.896, 0.917] & 0.896, 0.917 & 0.200 & no \\
 & \texttt{counts\_only} & 100 & 2{,}925 & 0.435 [0.417, 0.453] & 0.417, 0.453 & 0.020 & no \\
Windows & \texttt{no\_time} & 10 & 2{,}700 & 0.902 [0.890, 0.913] & 0.891, 0.913 & 0.200 & no \\
 & \texttt{no\_time} & 100 & 2{,}700 & 0.421 [0.403, 0.440] & 0.408, 0.442 & 0.020 & no \\
 & \texttt{counts\_only} & 10 & 2{,}700 & 0.902 [0.890, 0.913] & 0.890, 0.913 & 0.200 & no \\
 & \texttt{counts\_only} & 100 & 2{,}700 & 0.421 [0.403, 0.440] & 0.403, 0.440 & 0.020 & no \\
\midrule
\multicolumn{8}{l}{Recommendation from the two macOS cells: none} \\
\bottomrule
\end{tabular}
\end{sidewaystable}

% Generated by manuscript/tools/results_tables.py from the result files; do not edit by hand.
\begin{figure}[t]
\centering
\pgfplotslegendfromname{fig:closed-set:legend}\par\medskip
\begin{tikzpicture}
\begin{groupplot}[group style={group size=3 by 1, horizontal sep=0.45cm, yticklabels at=edge left, ylabels at=edge left}, fig panel, xmode=log, ymode=log, xtick={10,100,1000}, xmin=6, xmax=1700, ymin=0.0005, ymax=1.5, scale only axis, width=3.3cm, height=4.5cm, xlabel={$K$}, ylabel={Accuracy}]
\nextgroupplot[title={macOS, live conversion off}, legend to name=fig:closed-set:legend, legend columns=4]
\addlegendimage{fig set1, fig point}
\addlegendentry{\texttt{all}}
\addlegendimage{fig set2, fig point}
\addlegendentry{\texttt{no\_runs}}
\addlegendimage{fig set3, fig point}
\addlegendentry{\texttt{bin1s}}
\addlegendimage{fig set4, fig point}
\addlegendentry{\texttt{bin1s\_count}}
\addlegendimage{fig set5, fig point}
\addlegendentry{\texttt{no\_time}}
\addlegendimage{fig set6, fig point}
\addlegendentry{\texttt{counts\_only}}
\addlegendimage{fig chance}
\addlegendentry{$1/K$}
\addlegendimage{fig margin}
\addlegendentry{$2/K$}
\addplot[forget plot, fig chance, domain=7:1500] {1/x};
\addplot[forget plot, fig margin, domain=7:1500] {2/x};
\addplot[forget plot, fig interval, fig set1] table[row sep=\\] {x y\\7.94328 0.897056\\7.94328 0.91882\\7.94328 nan\\79.4328 0.41475\\79.4328 0.452106\\79.4328 nan\\794.328 0.0752358\\794.328 0.0963132\\794.328 nan\\};
\addplot[forget plot, fig set1, fig point] table[row sep=\\] {x y\\7.94328 0.908519\\79.4328 0.433333\\794.328 0.0851852\\};
\addplot[forget plot, fig interval, fig set2] table[row sep=\\] {x y\\8.70964 0.897056\\8.70964 0.91882\\8.70964 nan\\87.0964 0.41475\\87.0964 0.452106\\87.0964 nan\\870.964 0.0752358\\870.964 0.0963132\\870.964 nan\\};
\addplot[forget plot, fig set2, fig point] table[row sep=\\] {x y\\8.70964 0.908519\\87.0964 0.433333\\870.964 0.0851852\\};
\addplot[forget plot, fig interval, fig set3] table[row sep=\\] {x y\\9.54993 0.893163\\9.54993 0.915316\\9.54993 nan\\95.4993 0.41475\\95.4993 0.452106\\95.4993 nan\\954.993 0.0752358\\954.993 0.0963132\\954.993 nan\\};
\addplot[forget plot, fig set3, fig point] table[row sep=\\] {x y\\9.54993 0.904815\\95.4993 0.433333\\954.993 0.0851852\\};
\addplot[forget plot, fig interval, fig set4] table[row sep=\\] {x y\\10.4713 0.893163\\10.4713 0.915316\\10.4713 nan\\104.713 0.41475\\104.713 0.452106\\104.713 nan\\1047.13 0.0752358\\1047.13 0.0963132\\1047.13 nan\\};
\addplot[forget plot, fig set4, fig point] table[row sep=\\] {x y\\10.4713 0.904815\\104.713 0.433333\\1047.13 0.0851852\\};
\addplot[forget plot, fig interval, fig set5] table[row sep=\\] {x y\\11.4815 0.890052\\11.4815 0.91251\\11.4815 nan\\114.815 0.40298\\114.815 0.440206\\114.815 nan\\1148.15 0.0710503\\1148.15 0.0916225\\1148.15 nan\\};
\addplot[forget plot, fig set5, fig point] table[row sep=\\] {x y\\11.4815 0.901852\\114.815 0.421481\\1148.15 0.0807407\\};
\addplot[forget plot, fig interval, fig set6] table[row sep=\\] {x y\\12.5893 0.890052\\12.5893 0.91251\\12.5893 nan\\125.893 0.40298\\125.893 0.440206\\125.893 nan\\1258.93 0.064095\\1258.93 0.083784\\1258.93 nan\\};
\addplot[forget plot, fig set6, fig point] table[row sep=\\] {x y\\12.5893 0.901852\\125.893 0.421481\\1258.93 0.0733333\\};
\nextgroupplot[title={macOS, live conversion on}]
\addplot[forget plot, fig chance, domain=7:1500] {1/x};
\addplot[forget plot, fig margin, domain=7:1500] {2/x};
\addplot[forget plot, fig interval, fig set1] table[row sep=\\] {x y\\7.94328 0.895585\\7.94328 0.916682\\7.94328 nan\\79.4328 0.421082\\79.4328 0.457027\\79.4328 nan\\794.328 0.0742589\\794.328 0.0943543\\794.328 nan\\};
\addplot[forget plot, fig set1, fig point] table[row sep=\\] {x y\\7.94328 0.906667\\79.4328 0.438974\\794.328 0.0837607\\};
\addplot[forget plot, fig interval, fig set2] table[row sep=\\] {x y\\8.70964 0.895585\\8.70964 0.916682\\8.70964 nan\\87.0964 0.421082\\87.0964 0.457027\\87.0964 nan\\870.964 0.0742589\\870.964 0.0943543\\870.964 nan\\};
\addplot[forget plot, fig set2, fig point] table[row sep=\\] {x y\\8.70964 0.906667\\87.0964 0.438974\\870.964 0.0837607\\};
\addplot[forget plot, fig interval, fig set3] table[row sep=\\] {x y\\9.54993 0.895585\\9.54993 0.916682\\9.54993 nan\\95.4993 0.417343\\95.4993 0.453254\\95.4993 nan\\954.993 0.0742589\\954.993 0.0943543\\954.993 nan\\};
\addplot[forget plot, fig set3, fig point] table[row sep=\\] {x y\\9.54993 0.906667\\95.4993 0.435214\\954.993 0.0837607\\};
\addplot[forget plot, fig interval, fig set4] table[row sep=\\] {x y\\10.4713 0.895585\\10.4713 0.916682\\10.4713 nan\\104.713 0.417343\\104.713 0.453254\\104.713 nan\\1047.13 0.0742589\\1047.13 0.0943543\\1047.13 nan\\};
\addplot[forget plot, fig set4, fig point] table[row sep=\\] {x y\\10.4713 0.906667\\104.713 0.435214\\1047.13 0.0837607\\};
\addplot[forget plot, fig interval, fig set5] table[row sep=\\] {x y\\11.4815 0.895585\\11.4815 0.916682\\11.4815 nan\\114.815 0.417343\\114.815 0.453254\\114.815 nan\\1148.15 0.071355\\1148.15 0.0911124\\1148.15 nan\\};
\addplot[forget plot, fig set5, fig point] table[row sep=\\] {x y\\11.4815 0.906667\\114.815 0.435214\\1148.15 0.0806838\\};
\addplot[forget plot, fig interval, fig set6] table[row sep=\\] {x y\\12.5893 0.895585\\12.5893 0.916682\\12.5893 nan\\125.893 0.417343\\125.893 0.453254\\125.893 nan\\1258.93 0.067168\\1258.93 0.0864222\\1258.93 nan\\};
\addplot[forget plot, fig set6, fig point] table[row sep=\\] {x y\\12.5893 0.906667\\125.893 0.435214\\1258.93 0.0762393\\};
\nextgroupplot[title={Windows}]
\addplot[forget plot, fig chance, domain=7:1500] {1/x};
\addplot[forget plot, fig margin, domain=7:1500] {2/x};
\addplot[forget plot, fig interval, fig set1] table[row sep=\\] {x y\\7.94328 0.897056\\7.94328 0.91882\\7.94328 nan\\79.4328 0.414382\\79.4328 0.451735\\79.4328 nan\\794.328 0.0752358\\794.328 0.0963132\\794.328 nan\\};
\addplot[forget plot, fig set1, fig point] table[row sep=\\] {x y\\7.94328 0.908519\\79.4328 0.432963\\794.328 0.0851852\\};
\addplot[forget plot, fig interval, fig set2] table[row sep=\\] {x y\\8.70964 0.897056\\8.70964 0.91882\\8.70964 nan\\87.0964 0.414382\\87.0964 0.451735\\87.0964 nan\\870.964 0.0752358\\870.964 0.0963132\\870.964 nan\\};
\addplot[forget plot, fig set2, fig point] table[row sep=\\] {x y\\8.70964 0.908519\\87.0964 0.432963\\870.964 0.0851852\\};
\addplot[forget plot, fig interval, fig set3] table[row sep=\\] {x y\\9.54993 0.893163\\9.54993 0.915316\\9.54993 nan\\95.4993 0.414382\\95.4993 0.451735\\95.4993 nan\\954.993 0.0752358\\954.993 0.0963132\\954.993 nan\\};
\addplot[forget plot, fig set3, fig point] table[row sep=\\] {x y\\9.54993 0.904815\\95.4993 0.432963\\954.993 0.0851852\\};
\addplot[forget plot, fig interval, fig set4] table[row sep=\\] {x y\\10.4713 0.893163\\10.4713 0.915316\\10.4713 nan\\104.713 0.414382\\104.713 0.451735\\104.713 nan\\1047.13 0.0752358\\1047.13 0.0963132\\1047.13 nan\\};
\addplot[forget plot, fig set4, fig point] table[row sep=\\] {x y\\10.4713 0.904815\\104.713 0.432963\\1047.13 0.0851852\\};
\addplot[forget plot, fig interval, fig set5] table[row sep=\\] {x y\\11.4815 0.890052\\11.4815 0.91251\\11.4815 nan\\114.815 0.40298\\114.815 0.440206\\114.815 nan\\1148.15 0.0710503\\1148.15 0.0916225\\1148.15 nan\\};
\addplot[forget plot, fig set5, fig point] table[row sep=\\] {x y\\11.4815 0.901852\\114.815 0.421481\\1148.15 0.0807407\\};
\addplot[forget plot, fig interval, fig set6] table[row sep=\\] {x y\\12.5893 0.890052\\12.5893 0.91251\\12.5893 nan\\125.893 0.40298\\125.893 0.440206\\125.893 nan\\1258.93 0.064095\\1258.93 0.083784\\1258.93 nan\\};
\addplot[forget plot, fig set6, fig point] table[row sep=\\] {x y\\12.5893 0.901852\\125.893 0.421481\\1258.93 0.0733333\\};
\end{groupplot}
\end{tikzpicture}
\caption{Closed-set identification accuracy against the number of candidates $K$ for the six feature sets, per cell, with 95\% Wilson intervals; the accuracy of a set is the maximum over the set and the sets derivable from it. Solid line: the chance level $1/K$; dashed line: the margin $2/K$ of the near-chance rule, which judges \texttt{no\_time} and \texttt{counts\_only} at $K = 10$ and $100$; the sets that keep time are estimates only. The sets are set apart along $K$ for legibility; every point lies at $K = 10$, $100$ or $1{,}000$. Both axes are logarithmic. Table~\ref{tab:closed-set} gives the judgement. The figure is drawn from the result files by code written and reviewed with an AI tool (Claude) and checked by its tests; see the AI declaration}
\label{fig:closed-set}
\end{figure}

% Generated by manuscript/tools/results_tables.py from the result files; do not edit by hand.
\begin{sidewaystable}[p]
{\footnotesize
\caption{Open-set reconstruction per cell and attacker: character error rate (CER) of the attack and of the prior-only baseline, the gain (baseline CER minus attack CER) and the gain over the shuffled control (the CER of another trial's decoding minus the attack's), each with a block-clustered 95\% percentile interval. K1 holds out whole works, K1+ splits the corpus by phrase; the baseline knows the number of segments or the true length. Undecodable: trials for which no string fits; --: no value}
\label{tab:open-set}
\begin{tabular}{L{0.14}L{0.11}L{0.10}L{0.06}L{0.13}L{0.13}L{0.12}L{0.11}L{0.10}}
\toprule
Cell & Attacker & Length known & Trials & CER & Baseline CER & Gain & Gain over shuffled & Undecod\-able \\
\midrule
macOS, live conversion off & K0 & number of segments & 679 & 0.941 [0.937, 0.945] & 1.338 [1.281, 1.398] & 0.397 [0.337, 0.461] & 0.155 [0.124, 0.188] & 0 \\
 &  & true length & 679 & 0.941 [0.937, 0.945] & 0.968 [0.964, 0.971] & 0.027 [0.021, 0.032] & 0.155 [0.124, 0.188] & 0 \\
 & K1 & number of segments & 679 & 0.910 [0.903, 0.917] & 1.348 [1.282, 1.415] & 0.438 [0.371, 0.507] & 0.312 [0.260, 0.367] & 0 \\
 &  & true length & 679 & 0.910 [0.903, 0.917] & 0.961 [0.957, 0.965] & 0.051 [0.042, 0.059] & 0.312 [0.260, 0.367] & 0 \\
 & K1+ & number of segments & 544 & 0.913 [0.905, 0.922] & 1.353 [1.290, 1.422] & 0.439 [0.379, 0.508] & 0.334 [0.272, 0.401] & 0 \\
 &  & true length & 544 & 0.913 [0.905, 0.922] & 0.963 [0.958, 0.967] & 0.049 [0.041, 0.059] & 0.334 [0.272, 0.401] & 0 \\
 & K0 (split by phrase) & number of segments & 544 & 0.940 [0.935, 0.946] & 1.380 [1.319, 1.447] & 0.439 [0.377, 0.506] & 0.174 [0.134, 0.218] & 0 \\
 &  & true length & 544 & 0.940 [0.935, 0.946] & 0.969 [0.965, 0.972] & 0.028 [0.022, 0.035] & 0.174 [0.134, 0.218] & 0 \\
macOS, live conversion on & K0 & number of segments & 727 & 0.942 [0.938, 0.946] & 1.351 [1.292, 1.411] & 0.409 [0.348, 0.468] & 0.131 [0.105, 0.160] & 0 \\
 &  & true length & 727 & 0.942 [0.938, 0.946] & 0.967 [0.964, 0.970] & 0.026 [0.020, 0.031] & 0.131 [0.105, 0.160] & 0 \\
 & K1 & number of segments & 727 & 0.911 [0.903, 0.919] & 1.329 [1.273, 1.388] & 0.418 [0.359, 0.480] & 0.276 [0.232, 0.321] & 0 \\
 &  & true length & 727 & 0.911 [0.903, 0.919] & 0.964 [0.960, 0.968] & 0.053 [0.044, 0.062] & 0.276 [0.232, 0.321] & 0 \\
 & K1+ & number of segments & 580 & 0.916 [0.909, 0.924] & 1.329 [1.270, 1.394] & 0.413 [0.356, 0.471] & 0.290 [0.236, 0.349] & 0 \\
 &  & true length & 580 & 0.916 [0.909, 0.924] & 0.969 [0.965, 0.973] & 0.053 [0.044, 0.062] & 0.290 [0.236, 0.349] & 0 \\
 & K0 (split by phrase) & number of segments & 580 & 0.942 [0.937, 0.947] & 1.349 [1.292, 1.412] & 0.407 [0.349, 0.471] & 0.149 [0.116, 0.186] & 0 \\
 &  & true length & 580 & 0.942 [0.937, 0.947] & 0.967 [0.963, 0.971] & 0.025 [0.019, 0.030] & 0.149 [0.116, 0.186] & 0 \\
Windows & K0 & number of segments & 679 & 0.941 [0.937, 0.945] & 1.338 [1.281, 1.398] & 0.397 [0.337, 0.461] & 0.155 [0.124, 0.188] & 0 \\
 &  & true length & 679 & 0.941 [0.937, 0.945] & 0.968 [0.964, 0.971] & 0.027 [0.021, 0.032] & 0.155 [0.124, 0.188] & 0 \\
 & K1 & number of segments & 679 & 0.910 [0.903, 0.917] & 1.348 [1.282, 1.415] & 0.438 [0.371, 0.507] & 0.312 [0.260, 0.367] & 0 \\
 &  & true length & 679 & 0.910 [0.903, 0.917] & 0.961 [0.957, 0.965] & 0.051 [0.042, 0.059] & 0.312 [0.260, 0.367] & 0 \\
 & K1+ & number of segments & 544 & 0.913 [0.905, 0.922] & 1.353 [1.290, 1.422] & 0.439 [0.379, 0.508] & 0.334 [0.272, 0.401] & 0 \\
 &  & true length & 544 & 0.913 [0.905, 0.922] & 0.963 [0.958, 0.967] & 0.049 [0.041, 0.059] & 0.334 [0.272, 0.401] & 0 \\
 & K0 (split by phrase) & number of segments & 544 & 0.940 [0.935, 0.946] & 1.380 [1.319, 1.447] & 0.439 [0.377, 0.506] & 0.174 [0.134, 0.218] & 0 \\
 &  & true length & 544 & 0.940 [0.935, 0.946] & 0.969 [0.965, 0.972] & 0.028 [0.022, 0.035] & 0.174 [0.134, 0.218] & 0 \\
\bottomrule
\end{tabular}
}
\end{sidewaystable}


\begin{table}[t]
\caption{The near-chance rule: probability that a feature set is judged near chance when its true accuracy is $c/K$,
at the planned numbers of trials}
\label{tab:nearchance}
\begin{tabular}{rrrrrr}
\toprule
Trials & $K$ & $c=1$ & $c=1.5$ & $c=2$ & $c=3$ \\
\midrule
2{,}925 & 10 & 1.000 & 1.000 & 0.024 & 0.000 \\
2{,}925 & 100 & 0.994 & 0.487 & 0.020 & 0.000 \\
2{,}925 & 1{,}000 & 0.210 & 0.067 & 0.020 & 0.001 \\
2{,}700 & 10 & 1.000 & 1.000 & 0.025 & 0.000 \\
2{,}700 & 100 & 0.989 & 0.447 & 0.019 & 0.000 \\
2{,}700 & 1{,}000 & 0.067 & 0.017 & 0.004 & 0.000 \\
\bottomrule
\end{tabular}
\footnotetext{Binomial probability that the upper end of the Wilson interval is at most $2/K$, for the accuracy of
one feature set over the planned numbers of trials (2{,}925: live conversion on; 2{,}700: live conversion off and
Windows), as registered. For a set whose accuracy is the maximum over coarser sets the probability is at most this
value. If blocks are unscored, the rule is applied to the scored trials. The judgement of each cell is given in
Table~\ref{tab:closed-set}.}
\end{table}
